% Written by Anders Sjoqvist and Ulf Lundstrom, 2009
% The main sources are: tinyKACTL, Beta and Wikipedia

\chapter{Mathematics}

\section{Equations}
\[ax^2+bx+c=0 \Rightarrow x = \frac{-b\pm\sqrt{b^2-4ac}}{2a}\]

The extremum is given by $x = -b/2a$.

\[\begin{aligned}ax+by=e\\cx+dy=f\end{aligned}
\Rightarrow
\begin{aligned}x=\dfrac{ed-bf}{ad-bc}\\y=\dfrac{af-ec}{ad-bc}\end{aligned}\]

In general, given an equation $Ax = b$, the solution to a variable $x_i$ is given by
\[x_i = \frac{\det A_i'}{\det A} \]
where $A_i'$ is $A$ with the $i$'th column replaced by $b$.

\section{Recurrences}
If $a_n = c_1 a_{n-1} + \dots + c_k a_{n-k}$, and $r_1, \dots, r_k$ are distinct roots of $x^k - c_1 x^{k-1} - \dots - c_k$, there are $d_1, \dots, d_k$ s.t.
\[a_n = d_1r_1^n + \dots + d_kr_k^n. \]
Non-distinct roots $r$ become polynomial factors, e.g. $a_n = (d_1n + d_2)r^n$.

\section{Trigonometry}
\begin{align*}
	a\cos x+b\sin x&=r\cos(x-\phi)\\
	a\sin x+b\cos x&=r\sin(x+\phi)
\end{align*}
where $r=\sqrt{a^2+b^2}, \phi=\operatorname{atan2}(b,a)$.

\section{Geometry}

\subsection{Triangles}
Side lengths: $a,b,c$\\
Semiperimeter: $p=\dfrac{a+b+c}{2}$\\
Area: $A=\sqrt{p(p-a)(p-b)(p-c)}$\\
Circumradius: $R=\dfrac{abc}{4A}$\\
Inradius: $r=\dfrac{A}{p}$\\
Length of median (divides triangle into two equal-area triangles): $m_a=\tfrac{1}{2}\sqrt{2b^2+2c^2-a^2}$\\
Length of bisector (divides angles in two): $s_a=\sqrt{bc\left[1-\left(\dfrac{a}{b+c}\right)^2\right]}$\\
Law of sines: $\dfrac{\sin\alpha}{a}=\dfrac{\sin\beta}{b}=\dfrac{\sin\gamma}{c}=\dfrac{1}{2R}$\\
Law of cosines: $a^2=b^2+c^2-2bc\cos\alpha$\\
Law of tangents: $\dfrac{a+b}{a-b}=\dfrac{\tan\dfrac{\alpha+\beta}{2}}{\tan\dfrac{\alpha-\beta}{2}}$\\

\subsection{Quadrilaterals}
With side lengths $a,b,c,d$, diagonals $e, f$, diagonals angle $\theta$, area $A$ and
magic flux $F=b^2+d^2-a^2-c^2$:

\[ 4A = 2ef \cdot \sin\theta = F\tan\theta = \sqrt{4e^2f^2-F^2} \]

For cyclic quadrilaterals the sum of opposite angles is $180^\circ$,
$ef = ac + bd$, and $A = \sqrt{(p-a)(p-b)(p-c)(p-d)}$. \\

\subsection{Pick's Theorem}
For a simple polygon with integer coordinates:
\[ A = I + \frac{B}{2} - 1 \iff I = \frac{2A - B + 2}{2} \]

$I$: Number of integer lattice points strictly \textbf{inside} the polygon. \\
$B$: Number of integer lattice points on the \textbf{boundary} (including vertices): \\
$A$: Area of the polygon. \\
\[ B = \sum_{i=0}^{n-1} \gcd(|x_i - x_{(i+1)\%n}|, |y_i - y_{(i+1)\%n}|) \]

\textbf{Implementation Notes:} \\
\textbf{Avoid floats:} Use \texttt{polygonArea2} directly to get $2A = |\texttt{polygonArea2(pts)}|$, then compute $I = (2A - B + 2) / 2$ via integer arithmetic. \\
\textbf{Orientation:} Wrap \texttt{polygonArea2} with \texttt{abs()} since clockwise vertex ordering produces a negative signed area. \\
\textbf{Overflow:} Cast coordinates to \texttt{Point<long long>}; cross products overflow 32-bit signed integers when coordinates $\ge 10^5$. \\
\textbf{Ordering:} Vertices must appear in sequential order around the polygon boundary (e.g., from input or output of \texttt{convexHull}). \\

\section{Sums}
\[ c^a + c^{a+1} + \dots + c^{b} = \frac{c^{b+1} - c^a}{c-1}, c \neq 1 \]
\begin{align*}
	1 + 2 + 3 + \dots + n &= \frac{n(n+1)}{2} \\
	1^2 + 2^2 + 3^2 + \dots + n^2 &= \frac{n(2n+1)(n+1)}{6} \\
	1^3 + 2^3 + 3^3 + \dots + n^3 &= \frac{n^2(n+1)^2}{4} \\
	1^4 + 2^4 + 3^4 + \dots + n^4 &= \frac{n(n+1)(2n+1)(3n^2 + 3n - 1)}{30} \\
\end{align*}

\section{Stars and Bars}
Number of ways to distribute $n$ identical items into $k$ distinct bins (or number of integer solutions to $x_1 + x_2 + \dots + x_k = n$):
\textbf{Non-negative integers} ($x_i \ge 0$, empty bins allowed):
\[ \binom{n + k - 1}{k - 1} = \binom{n + k - 1}{n} \]
\textbf{Positive integers} ($x_i \ge 1$, each bin has at least 1 item, $n \ge k$):
\[ \binom{n - 1}{k - 1} \]
\textbf{Arbitrary lower bounds} ($x_i \ge c_i$): Let $n' = n - \sum_{i=1}^k c_i$.
If $n' \ge 0$:
\[ \binom{n' + k - 1}{k - 1} \]
\textbf{Upper bounds} ($x_i \le u_i$): Apply the Principle of Inclusion-Exclusion by negating constraints ($x_i \ge u_i + 1$).

\subsection{Geometric Series Modulo Arbitrary M}
To compute $S(n) = \sum_{i=0}^{n} a^i = 1 + a + a^2 + \dots + a^n \pmod M$ in $\mathcal{O}(\log^2 n)$ (or $\mathcal{O}(\log n)$) without requiring modular inverse (valid for composite $M$):
If $n = 0$: $S(0) = 1$. \\
If $n$ is odd ($n+1$ terms is even): \\
\[ S(n) = \left(1 + a^{\frac{n+1}{2}}\right) S\left(\frac{n-1}{2}\right) \pmod M \]
If $n$ is even ($n+1$ terms is odd): \\
\[ S(n) = 1 + a \cdot S(n-1) \pmod M \]

\begin{lstlisting}[language=C++]
ll geomSeries(ll a, ll n, ll m) {
	if (n == 0) return 1 % m;
	if (n % 2) {
		ll half = geomSeries(a, (n - 1) / 2, m);
		ll p = modpow(a, (n + 1) / 2, m);
		return (__int128)(1 + p) * half % m;
	}
	return (1 + (__int128)(a % m) * geomSeries(a, n - 1, m)) % m;
}
\end{lstlisting}

\section{Probability theory}
Let $X$ be a discrete random variable with probability $p_X(x)$ of assuming the value $x$. It will then have an expected value (mean) $\mu=\mathbb{E}(X)=\sum_xxp_X(x)$ and variance $\sigma^2=V(X)=\mathbb{E}(X^2)-(\mathbb{E}(X))^2=\sum_x(x-\mathbb{E}(X))^2p_X(x)$ where $\sigma$ is the standard deviation. If $X$ is instead continuous it will have a probability density function $f_X(x)$ and the sums above will instead be integrals with $p_X(x)$ replaced by $f_X(x)$.

Expectation is linear:
\[\mathbb{E}(aX+bY) = a\mathbb{E}(X)+b\mathbb{E}(Y)\]
For independent $X$ and $Y$, \[V(aX+bY) = a^2V(X)+b^2V(Y).\]

\subsection{Discrete distributions}

\subsubsection{Binomial distribution}
The number of successes in $n$ independent yes/no experiments, each which yields success with probability $p$ is $\textrm{Bin}(n,p),\,n=1,2,\dots,\, 0\leq p\leq1$.
\[p(k)=\binom{n}{k}p^k(1-p)^{n-k}\]
\[\mu = np,\,\sigma^2=np(1-p)\]
$\textrm{Bin}(n,p)$ is approximately $\textrm{Po}(np)$ for small $p$.

\subsubsection{First success distribution}
The number of trials needed to get the first success in independent yes/no experiments, each which yields success with probability $p$ is $\textrm{Fs}(p),\,0\leq p\leq1$.
\[p(k)=p(1-p)^{k-1},\,k=1,2,\dots\]
\[\mu = \frac1p,\,\sigma^2=\frac{1-p}{p^2}\]
